% ==========================================
% 1. 论文元信息
% ==========================================
\documentclass[fontset=none, UTF8]{ctexart}
\def\MYTITLEA{请输入你的论文标题}
\def\MYTITLEB{如果太长了可以换行}
\def\MYNAME{738NGX}
\def\MYID{2022114514}
\def\MYCOLLEGE{计算机与人工智能学院}
\def\MYMAJOR{计算机科学与技术}
\def\MYADVISOR{天王寺璃奈}
\def\MYDATE{ \quad 年  \quad \quad 月}
% ==========================================
% 1. 基础宏包与数学支持
% ==========================================
\usepackage{amsmath, amsfonts, amssymb}
\usepackage{graphicx}
\usepackage{xstring}
\usepackage{amssymb} % 用于实现 \checkmark 钩号

% 备份原始命令
\let\oldincludegraphics\includegraphics

% 重定义 \includegraphics 以自动剥离 .png 后缀，如不需要可以注释或删除
\renewcommand{\includegraphics}[2][]{%
    \IfEndWith{#2}{.png}{%
        \StrBefore{#2}{.png}[\filename]% 如果以 .png 结尾，截取前面部分
        \oldincludegraphics[#1]{\filename}%
    }{%
        \oldincludegraphics[#1]{#2}% 如果不是，按原样处理
    }%
}

% 适配某些版本的pandoc转换可能导致的错误
\makeatletter
\newsavebox\pandoc@box
\newcommand*\pandocbounded[1]{%
  \sbox\pandoc@box{#1}%
  \ifdim\wd\pandoc@box>\linewidth
    \resizebox{\linewidth}{!}{\usebox\pandoc@box}%
  \else
    \usebox\pandoc@box
  \fi
}
\makeatother

% 确保 PDF 优先级最高
\DeclareGraphicsExtensions{.pdf,.png,.jpg}
\usepackage{float}
\usepackage{pdfpages}
\usepackage{setspace}
\usepackage{microtype}
\usepackage{booktabs}   % 常用表格线宏包
\usepackage{longtable}  % 支持跨页长表格
\usepackage{array}      % 增强表格功能

% ==========================================
% 2. 字体与中文设置
% ==========================================
\usepackage{xeCJK}

% 基础字体设置
\setmainfont{Times New Roman}
\setCJKmainfont{SIMSUN.TTC}[AutoFakeBold=2.5]
\setsansfont{Arial}
\setCJKsansfont{SIMHEI.TTF}[AutoFakeBold=2.5]
\setCJKmonofont{SIMHEI.TTF}
\setCJKfamilyfont{zhkai}{SIMKAI.TTF}[AutoFakeBold=2.5]
\newcommand{\kaishu}{\CJKfamily{zhkai}}

% ==========================================
% 3. 页面边距与页码格式
% ==========================================
\usepackage[a4paper, top=2.5cm, bottom=2.5cm, left=3cm, right=2cm, footskip=1.0cm]{geometry}
\usepackage{fancyhdr}

\fancypagestyle{plain}{
    \fancyhf{}
    \fancyfoot[C]{{\zihao{-5} 第 \thepage 页}}
    \renewcommand{\headrulewidth}{0pt}
}
\pagestyle{plain}

% ==========================================
% 4. 参考文献 (GB/T 7714-2015)
% ==========================================
\usepackage[
    backend=biber,
    style=gb7714-2015,
    gbalign=left,
    gbpub=false,
    autocite=superscript
]{biblatex}
\addbibresource{sample.bib}

% ==========================================
% 5. 图表格式与编号设置
% ==========================================
\usepackage{caption}
\usepackage{chngcntr}

% 图号格式：图 3-1
\counterwithin{figure}{section}
\renewcommand{\thefigure}{\thesection--\arabic{figure}}
\renewcommand{\figurename}{图}

% 强制图片位置与尺寸限制
\makeatletter
\def\maxwidth{\ifdim\Gin@nat@width>\linewidth\linewidth\else\Gin@nat@width\fi}
\makeatother
\setkeys{Gin}{width=\maxwidth, height=0.3\textheight, keepaspectratio}

\DeclareGraphicsRule{.png}{pdf}{.pdf}{} % 优先加载同名 pdf
\floatplacement{figure}{H}              % 强制 H 选项

% 自定义图表名格式
\DeclareCaptionFormat{custom}{\kaishu\zihao{5} #1#2#3}
\captionsetup[figure]{
    format=custom,
    labelsep=quad,
    position=bottom,
    font={small}
}
\counterwithin{table}{section}
\renewcommand{\thetable}{\thesection--\arabic{table}}
\renewcommand{\tablename}{表}
\captionsetup[table]{
    format=custom,
    labelsep=quad,
    position=top,
    font={small}
}

% ==========================================
% 6. 章节标题格式 (ctexset)
% ==========================================
\ctexset{
  section = {
    name = {, 、}, 
    number = \chinese{section},
    format = \zihao{3}\sffamily\centering,
    aftername = \hspace{0pt},
    beforeskip = 1ex,
    afterskip = 24pt
  },
  subsection = {
    name = {（, ）},
    number = \chinese{subsection},
    format = \zihao{4}\sffamily\raggedright,
    aftername = \hspace{0pt},
    beforeskip = 1ex,
    afterskip = 12pt
  },
  subsubsection = {
    name = {, },
    number = \arabic{subsubsection}．,
    format = \zihao{5}\sffamily\raggedright,
    aftername = \hspace{0pt},
    beforeskip = 1ex,
    afterskip = 1ex
  }
}

% ==========================================
% 7. 目录格式 (titletoc)
% ==========================================
\usepackage{titletoc}
\setcounter{tocdepth}{3}
\setcounter{secnumdepth}{3}
\renewcommand{\contentsname}{\hfill\zihao{3}\sffamily 目\quad 录\hfill}

\titlecontents{section}
    [0em] {\addvspace{6pt}\zihao{4}\sffamily}
    {\makebox[3em][l]{\thecontentslabel}} {} 
    {\titlerule*[0.5pc]{.}\contentspage}

\titlecontents{subsection}
    [2em] {\addvspace{2pt}\zihao{-4}\sffamily}
    {\makebox[2.5em][l]{\thecontentslabel}} {}
    {\titlerule*[0.5pc]{.}\contentspage}

\titlecontents{subsubsection}
    [4em] {\addvspace{2pt}\zihao{-4}\sffamily}
    {\makebox[2em][l]{\thecontentslabel}} {}
    {\titlerule*[0.5pc]{.}\contentspage}

% ==========================================
% 8. 其他兼容性修正与超链接
% ==========================================
% 解决 Pandoc 的 \tightlist 报错
\providecommand{\tightlist}{%
  \setlength{\itemsep}{0pt}\setlength{\parskip}{0pt}}
% 超链接设置（最后加载）
\usepackage[unicode=true]{hyperref}
\hypersetup{hidelinks}
% 正文基础字号
\renewcommand{\normalsize}{\zihao{5}}


% ==========================================
% 正文开始
% ==========================================
\begin{document}

% 1. 导入封面
%\includepdf[pages={1,2}]{cover.pdf}
% ==========================================
% 封面与声明页
% ==========================================

\newgeometry{left=3cm, right=3cm, top=2.5cm, bottom=2.5cm}

\begin{titlepage}
    \pagestyle{empty} 
    
    {
        \centering
        \vspace*{0.5cm}
        {\zihao{-0} 上海财经大学 \par}
        \vspace{1.5cm}
        {\zihao{1} 毕\quad 业\quad 论\quad 文 \par}
        
        \vspace{2cm}

        % 题目部分：直接调用新定义的 \MYTITLEA
        {\zihao{2}
        \begin{tabular}{rc}
            \makebox[2em][s]{题\hfill 目} & \underline{\makebox[10em][c]{\sffamily \MYTITLEA}} \\
            & \underline{\makebox[10em][c]{\sffamily \MYTITLEB}} \\
        \end{tabular}}

        \vspace{3cm}

        % 个人信息栏
        {\zihao{4}
        \renewcommand{\arraystretch}{1.8} 
        \begin{tabular}{lc}
            \makebox[4em][s]{姓\hfill 名} & \underline{\makebox[10em][c]{\MYNAME}} \\
            \makebox[4em][s]{学\hfill 号} & \underline{\makebox[10em][c]{\MYID}} \\
            \makebox[4em][s]{学\hfill 院} & \underline{\makebox[10em][c]{\MYCOLLEGE}} \\
            \makebox[4em][s]{专\hfill 业} & \underline{\makebox[10em][c]{\MYMAJOR}} \\
            \makebox[4em][s]{指导教师} & \underline{\makebox[10em][c]{\MYADVISOR}} \\
        \end{tabular}}

        \vfill
        {\zihao{4} 定稿日期 \hspace{4em} \MYDATE \par}
        \vspace{1cm}
    }
    
    \clearpage

    % 声明页
    \begin{center}
        \vspace*{0.5cm}
        {\zihao{1} \bfseries\sffamily 声 \qquad 明 \par}
    \end{center}

    \vspace{1.5cm}

    \begin{spacing}{1.5} 
        \zihao{2}\kaishu\bfseries
        \indent 本人郑重声明所呈交的论文是我个人在指导老师的指导下进行的研究工作及取得的研究成果，不存在任何剽窃、抄袭他人学术成果的现象。本人确保在论文写作中严格遵守学术行为规范，未使用人工智能技术取代论文写作，符合学术诚信要求。
        
        \indent 我同意（ $\checkmark$ ）/ 不同意（ \qquad ）本论文作为学校的信息资料使用。
    \end{spacing}

    \vfill

    \begin{flushright}
        \zihao{2}\kaishu\bfseries
        \begin{tabular}{r c}
            论文作者（签名）： & \underline{\makebox[6em]{}} \\ [1.5cm]
            & \makebox[6em][s]{年 \hfill 月 \hfill 日} 
        \end{tabular}
    \end{flushright}
    \vspace{1cm}

\end{titlepage}

\restoregeometry

% 2. 摘要页初始化
\restoregeometry 
\newpage
\pagestyle{plain} 
\pagenumbering{Roman} 
\setcounter{page}{1}

% --- 中文摘要 ---
\addcontentsline{toc}{section}{\texorpdfstring{摘\quad 要}{摘要}}
\begin{center}
    {\zihao{3}\sffamily 论文标题} 
\end{center}
\vspace{2em}
\begin{center}
    {\zihao{3}\sffamily 摘\quad 要}
\end{center}

\vspace{1em}

\begin{spacing}{2.5}
{\zihao{4}
请在此处写入正文摘要。
}
\end{spacing}

\vspace{2em}
\noindent {\zihao{4}\sffamily 关键词：}{\zihao{4} 关键词A，关键词B，关键词C}

\newpage
% --- 英文摘要 ---
\begin{center}
{\zihao{3}\sffamily 
Title
}
\end{center}
\vspace{1.5em}
\begin{center}
{\zihao{3}\sffamily Abstract}
\end{center}
\addcontentsline{toc}{section}{Abstract}

\vspace{1em}

{\zihao{5}\sffamily
Please input english abstract here.
}

\vspace{2em}
\noindent {\zihao{5}\sffamily\bfseries Key words: }{\zihao{5}\sffamily 
Key A, Key B, Key C.
} 

\newpage
\tableofcontents 

\newpage
\pagenumbering{arabic} 
\setcounter{page}{1}

\begin{spacing}{1.5}

% --- 摘要页初始化 ---
\pagestyle{plain} 
\pagenumbering{Roman} 
\setcounter{page}{1}

% --- 中文摘要 ---
\addcontentsline{toc}{section}{\texorpdfstring{摘\quad 要}{摘要}}
\begin{center}{\zihao{3}\sffamily 请输入论文的中文标题，将在摘要页面展示} \end{center}
\vspace{2em}
\begin{center}{\zihao{3}\sffamily 摘\quad 要}\end{center}
\vspace{1em}
\begin{spacing}{2.5}{\zihao{4} 这里写入你的中文摘要内容。}\end{spacing}
\vspace{2em}
\noindent {\zihao{4}\sffamily 关键词：}{\zihao{4} 关键词A; 关键词B;
关键词C}

\clearpage

% --- 英文摘要 ---
\addcontentsline{toc}{section}{Abstract}
\begin{center}{\zihao{3}\sffamily Please input the English
title}\end{center}
\vspace{1.5em}\begin{center}{\zihao{3}\sffamily Abstract}\end{center}
\vspace{1em}
\begin{spacing}{1}{\zihao{5}\sffamily Enter your English abstract
here.}\end{spacing}
\vspace{2em}
\noindent {\zihao{5}\sffamily\bfseries Key words: }{\zihao{5}\sffamily Keyword
A; Keyword B; Keyword C}

% --- 目录 ---
\clearpage
\tableofcontents 

% --- 正文 ---
\clearpage
\pagenumbering{arabic} 
\setcounter{page}{1}
\begin{spacing}{1.5}
\section{正文一级标题}

\subsection{正文二级标题}

\subsubsection{正文三级标题}

这是一段用于测试文本样式的长正文文本，包含\textbf{粗体}、\emph{斜体}和\textbf{\emph{粗体+斜体}}：\textbf{LoveLive!}（日语：\textbf{ラブライブ！}），全称\textbf{\emph{LoveLive!
School idol project}}（日语：\textbf{ラブライブ！ School idol
project}），是由日本动画公司日昇、唱片公司Lantis，以及 ACG 月刊杂志
\emph{《电击G's magazine》} 在 2010
年共同合作推出的读者参与型混合媒体校园偶像企划，企划包括音乐作品、电视和广播节目、文字书籍、
漫画插画 、动画 、游戏、周边等。

这是一段用于测试列表样式的正文文本，包含无序列表和有序列表。

无序列表

\begin{itemize}
\tightlist
\item
  LoveLive!
\item
  BanGDream
\item
  THE IDOLM@STER
\end{itemize}

有序列表

\begin{enumerate}
\def\labelenumi{\arabic{enumi}.}
\tightlist
\item
  高坂穗乃果
\item
  户山香澄
\item
  天海春香
\end{enumerate}

这是一段用于测试公式样式的正文文本，包括行内公式
\(= LOVE\)、\(\neq ME\)、\(\fallingdotseq JOY\)。和独立公式： \[
\text{ユメ}+\text{ミライ}=\text{無限大}
\]
这是一段用于展示图片使用的正文文本。正确按本示例书写引用时图片按照格式要求会自动进行编号（见图
\ref{fig:sample-fig}）。图片的最大高度被限制为最大占用页面文字区域高度的
30\%，如有需要可自行进行调整。推荐使用pdf格式的图片来提升编译效率。模板包含了将同名png后缀自动修改为pdf的功能，这代表你可以在图片文件夹下同时放入两张同名的
pdf 和 png 文件来兼顾 markdown 显示和 latex 编译。

\begin{figure}
\centering
\includegraphics{./figures/sample.png}
\caption{示例图片}\label{fig:sample-fig}
\end{figure}

这是一段用于展示表格使用的正文文本。正确按本示例书写引用时表格按照格式要求会自动进行编号（见表
\ref{tbl:sample-tbl}）。

\begin{longtable}[]{@{}llllll@{}}
\caption{示例表格 \label{tbl:sample-tbl}}\tabularnewline
\toprule\noalign{}
MyGO!!!!! & 主唱 & 主音吉他 & 节奏吉他 & 贝斯 & 鼓手 \\
\midrule\noalign{}
\endfirsthead
\toprule\noalign{}
MyGO!!!!! & 主唱 & 主音吉他 & 节奏吉他 & 贝斯 & 鼓手 \\
\midrule\noalign{}
\endhead
\bottomrule\noalign{}
\endlastfoot
角色 & 高松灯 & 要乐奈 & 千早爱音 & 长崎素世 & 椎名立希 \\
声优 & 羊宫妃那 & 青木阳菜 & 立石凛 & 小日向美香 & 林鼓子 \\
\end{longtable}

这是一段引用了参考文献的正文文本\autocite{yuJapaneseAnimeIdol2021}。参考文献的角标具体样式可以自行调整。参考文献会在文末一并列出。
\end{spacing}

% --- 参考文献 ---
\clearpage
\printbibliography[heading=bibintoc, title={参考文献}]


\end{spacing}

\newpage
\printbibliography[heading=bibintoc, title={参考文献}]

\end{document}